\section{Introduction}
\label{sec:intro}

Large Language Models (LLMs) have revolutionized artificial intelligence (AI) by delivering unprecedented capabilities in natural language processing. These models are built on advanced deep learning architectures known as transformer networks, which excel at handling sequential data and understanding context \citep{vaswani2017attention,radford2018improving}. Trained on extensive and diverse datasets—including vast amounts of text from books, articles, and websites—LLMs are capable of generating human-like responses and performing complex language tasks. For example, Llama 3 was trained on 15 trillion tokens of data \citep{llama}. This training enables LLMs to comprehend and produce coherent, contextually relevant text, making them invaluable tools in a wide range of applications \citep{brown2020language}.

The rapid development of LLMs also signifies a transformative moment in the social sciences. In market research—a critical area in marketing science and operations management—understanding market demand for products or services is paramount. Traditional market research methods, such as surveys, focus groups, and conjoint analysis, often require significant resources and time. Gathering data from real subjects involves designing intricate questionnaires, recruiting participants, and analyzing responses, all of which can be both costly and logistically challenging \citep{hair2019essentials}. LLMs present a novel solution by enabling the generation of synthetic data that simulates real consumer behavior without the need for extensive fieldwork. This potential has sparked growing interest among researchers. Thus, a key advantage of using LLMs in market research is their ability to scale data generation: whereas traditional methods often limit the volume of data that can be feasibly collected due to constraints on time and resources, LLMs can rapidly generate large volumes of data, facilitating more comprehensive and granular analysis.
 
LLMs are trained on vast amounts of data from diverse sources across the internet, including reviews, purchase decisions, and even real survey data on products and services from various categories. The underlying assumption, expressed with some exaggeration, is that LLMs have been exposed to everything there is to see. Researchers hypothesize that the responses LLMs generate for market research surveys could, to some extent, reflect the types of responses that real consumers---represented in the training data---would have given to the same questions. These factors suggest that LLMs may become invaluable tools for gaining insights into consumer preferences, thanks to their ability to mimic or replicate human responses \citep{brand2023using}.
Recent studies provided encouraging evidence supporting this potential \citep{argyle2023out,chen2023emergence,horton2023large}. For instance, research showed that by fine-tuning a model \citep{gururangan2020don} or employing prompting techniques such as chain-of-thought (CoT) reasoning \citep{wei2022chain}, it is possible to generate simulated responses in market surveys that resemble those of real human subjects \citep{brand2023using,goli2024frontiers}. These explorations suggest the enormous potential of LLMs for advancing market research.

Nonetheless, recent literature also highlights several caveats. Despite some encouraging findings by \cite{goli2024frontiers}, the authors acknowledged that even the most up-to-date prompt engineering methods and advanced LLMs, such as OpenAI's Generative Pre-trained Transformer (GPT), cannot completely eliminate discrepancies between human and LLM responses in market surveys. For example, they observe that GPT-4 exhibits a pronounced level of impatience compared to human participants. This led the authors to conclude that \textit{``directly eliciting preferences using LLMs can yield misleading results,''} suggesting that LLMs should be used more as tools to facilitate hypothesis exploration rather than as direct substitutes for human responses. Similarly, \cite{gui2023challenge} pointed out that LLMs tend to make unintended assumptions about unspecified details of the choice environment, further complicating their use in market research.

We acknowledge the significant potential of LLM-generated labels for market research, but we also believe that the aforementioned gap is likely to persist, even with the most advanced future developments in AI. Consider conjoint analysis, which is arguably the most fundamental methodology for understanding consumer preferences and decision-making processes in market research \citep{green1990conjoint}. Typically, the researcher seeks to understand how consumers value different features or attributes of a product or service (for details on the setup, please refer to Section \ref{sec:model}).
At its core, customers aim to maximize their benefits by selecting the optimal products based on their past experiences. However, we emphasize that LLMs do not possess real-life experiences. Additionally, consumer preferences evolve over time due to trends, technological advancements, and economic and cultural shifts \citep{solomon2020consumer}, changes that LLMs may not accurately capture.

The gap between an LLM's predictions and actual customer preferences can lead to misleading results if a researcher simply replaces real data with AI-generated data or na\"ively combines the two—a point underscored in \cite{goli2024frontiers} through careful empirical studies. This observation aligns with our findings. We conducted an empirical study on the conjoint analysis of COVID-19 vaccine preferences, pooling real data with LLM-generated responses based on the same set of product features and performing the conjoint analysis in a manner similar to \cite{brand2023using} (see Section \ref{sec:empirical-1} for details of the study).
The results are presented in Figure \ref{fig:llm-human}. 
Regardless of the version of the LLM used or whether state-of-the-art prompting techniques are applied, there is a significant gap between the ground-truth model coefficients and those estimated using pooled data. This gap does not diminish with the inclusion of more LLM-generated data; in most models, it actually worsens as more LLM-generated data is added. In contrast, as one would expect, the gap decreases when more real data is incorporated, eventually vanishing to zero, highlighting a stark contrast with the LLM-generated data.


\begin{figure}[!t]
    \centering
    % \includegraphics[width=100mm]{fig/Picture1.jpg}
    \includegraphics[width=0.85\textwidth]{fig/bias-naive.pdf}
    \caption{LLM-generated Data $\neq$ Human Data}
    \label{fig:llm-human}
    \begin{minipage}{\textwidth}
    \footnotesize
    \begin{tabular}{p{1cm} p{15cm}}
         \hspace{-1mm}\textbf{Notes}: & \setlength{\baselineskip}{1.2\baselineskip} The plot illustrates the percentage error of estimation, calculated as \(\big\lVert (\hat{\betavec} - \betavec^*)/{\betavec^*} \big\rVert_1\), where the ratio is computed component-wise. This error is based on pooling 50 real customer response data points with LLM-generated data of varying sizes. Each curve represents estimators obtained using different combinations of AI-generated or real labels. For further details, please refer to Section \ref{sec:empirical-1}. The horizontal axis indicates the quantity of data pooled into the model.
    \end{tabular}
    
    
    % \begin{itemize}[leftmargin=1cm]
    %     \item [\textbf{Notes}:] The plot shows the percentage error of estimation based on pooling the real customer response data and LLM-generated data, $\big\lVert {\big(\hat{\betavec} - \betavec^*\big)}/{\betavec^*}\big\rVert_1$, where the ratio is taken component-wise. Different curve corresponds to estimators obtained using different AI-generated or real labels. For more details, please refer
    %   to Section \ref{sec:empirical-1}. The horizontal axis gives the quantity of data pooled into the model.
    % \end{itemize}
  \end{minipage}
\end{figure}

Recognizing that AI-generated data, while informative, is inherently imperfect and likely to be misleading if used na\"ively, our goal is to develop a scheme for its efficient utilization. Although engineering techniques such as fine-tuning, CoT prompting, and in-context learning can potentially improve the quality of data from LLMs, our objective is to introduce a statistical approach that can be broadly applied independently of these methods and have theoretical guarantees. Our research question is then: \emph{Can we improve conjoint analysis by effectively integrating LLM-generated data with real data to produce more accurate estimators of customer preferences, compared to using human data, LLM-generated data, or a na\"ive combination of both?} Our method is inspired by the widely recognized transfer learning and knowledge distillation approach in the deep learning literature \citep{hinton2015distilling}. By formalizing and adapting their core principles to our context, we present a data-augmentation statistical approach for extracting value from LLM- and AI-generated conjoint data. Specifically, we make two key contributions.

\vspace{1mm}
\noindent
\revise{\underline{\textit{Novel Theoretical Framework for Data Augmentation.}} 
We develop a general framework for statistical data augmentation using generative models, such as LLMs, to capture human preferences in conjoint analysis, and propose the AI-Augmented Estimator
(AAE). The framework is broadly applicable to hybrid datasets where true human labels (ground truth) are available for only a subset of data points, while (AI-)generated labels can be created for all observations. These AI-generated labels may be produced by a traditional machine learning model or a generative model such as an LLM. Although our framework is compatible with any source of surrogate labels, results presented in Section \ref{sec:empirical-1} and Appendix \ref{app:comparison} show that the estimator performs best when the AI-generated data are more informative. This finding suggests that, at the time of writing, the most effective application of our framework is to augment data generated by LLMs.}

\revise{
Conceptually, our approach leverages the joint distribution among features, surrogate labels, and true labels within the subset of data where both types of labels are observed. This joint structure allows us to learn the distributional discrepancy between surrogate and true labels and subsequently regularize the surrogate labels so that they more closely approximate the true ones. By doing so, we effectively expand the labeled dataset, enabling more accurate estimation of human preferences. These principles give rise to theoretically grounded estimators: our estimator is consistent and asymptotically normal (Theorem \ref{thm:consistency}) and satisfies a finite-sample, high-probability performance bound (Theorem \ref{thm:aae-bound}). Furthermore, we demonstrate that our method outperforms approaches that rely solely on human data, solely on LLM-generated data, or on naïve combinations of the two.}

\vspace{1mm}
\noindent
\revise{
\underline{\textit{Large-scale 
 empirical validation.}} We conduct extensive empirical analyses across multiple conjoint settings to validate the effectiveness and robustness of our proposed AAE. In the first empirical study, based on a real-world COVID-19 vaccination dataset, AAE consistently reduces estimation error relative to naïve augmentation and AI-only baselines, even when GPT-generated labels are imperfect. Its regularization capability remains robust across different LLM versions and prompting strategies—basic, CoT, few-shot, and fine-tuned—with GPT-4o and CoT prompting achieving the greatest gains. The results confirm that combining LLM-generated and human data through a statistically grounded augmentation framework yields reliable, unbiased preference estimates while reducing survey data requirements by 24.9\% to 79.8\%. In Section 6, we extend these findings to a separate sports-car choice dataset and additional robustness checks. AAE again outperforms alternative estimators across all models and prompts, outperforming even fine-tuned GPT-4o baselines. Further comparisons show that AAE matches  prediction-powered inference (PPI), an alternative data augmentation strategy
 \cite[see the literature review and ][]{angelopoulos2023prediction, angelopoulos2023ppi++} in low-dimensional cases but surpasses it in high-dimensional conjoint features due to its more granular correction of AI-human discrepancies.
 % Finally, under model misspecification and nonparametric estimation settings, AAE continues to exhibit strong stability and accuracy. 
 Collectively, these empirical results demonstrate that AAE offers a scalable, bias-controlled, and cost-efficient solution for integrating LLM-generated data into real-world market research.}


We briefly discuss the organization of our paper. In Section \ref{sec:literature}, we present a literature review.
In Section \ref{sec:model}, we set up our model and formally state our data-augmentation approach. We discuss the theoretical properties of our proposed approach in Section \ref{sec:consistency-normality} and two empirical studies in Sections \ref{sec:empirical-1} and \ref{sec:empirical-2}. In Section \ref{sec:conclusion}, we conclude. 